% !TeX root = ../example.tex
\section{基础页面与图像}
% 本示例每节显式插入起始页：图片直接在本次汇报中选择。
% 实际使用时，可按需省略目录页或起始页，也可交换两者顺序。
\sectioncoverpage[assets/sigs_building]{本节介绍页面布局、公式排版、图像适配与图注。}
\note{本节先说明正文区域和字号，再演示公式、代码与图片。章节背景是模板已有示例素材。}
\subsection{基础页面}
% !TeX root = ../example.tex
\begin{frame}{先确认大纲，再制作汇报}
      \begin{description}[最终交付]
            \item[准备内容] 在根目录 \texttt{outline.md} 写草稿，素材放入 \texttt{materials/}。
            \item[确认大纲] agent 整理章节、每页主旨、正文与备注分工，交由你审阅确认。
            \item[制作汇报] 确认后生成正文与备注，检查页面并导出 PDF/PPTX。
            \item[最终交付] \texttt{build/main.pdf} 和带逐页备注的 \texttt{build/main.pptx}。
      \end{description}
      \medskip
      对 agent 说：根据大纲制作 PPT。
\end{frame}
 % 推荐用法：思路与素材交给 agent

% 红色参考线只显示在这一页，不改变其他页面的主题设置。
% Diagnostic instrumentation is local to this frame; the theme is unchanged.
\begingroup
  % Horizontal dimensions use mm; vertical dimensions use the named font's
  % baseline grid (header: tiny/Large/normalsize; footer: scriptsize).
\newcommand{\layoutGuideDistance}[2]{%
  \begingroup
    \pgfextractx{\dimen0}{\pgfpointdiff{\pgfpointanchor{#1}{center}}{\pgfpointanchor{#2}{center}}}%
    \pgfextracty{\dimen2}{\pgfpointdiff{\pgfpointanchor{#1}{center}}{\pgfpointanchor{#2}{center}}}%
    \pgfmathparse{veclen(\the\dimen0,\the\dimen2)/(1mm)}%
    \pgfmathprintnumber[fixed,precision=2]{\pgfmathresult}\,mm%
  \endgroup
}
% Vertical dimensions are expressed in the same grid unit used by the
% navigation and footer rows. Horizontal dimensions continue to use mm.
\newcommand{\layoutGuideBaselineDistance}[3][\campusfooterbaseline]{%
  \begingroup
    \pgfextracty{\dimen2}{\pgfpointdiff{\pgfpointanchor{#2}{center}}{\pgfpointanchor{#3}{center}}}%
    \pgfmathparse{abs(\the\dimen2)/(1mm) / ((\the#1)/(1mm))}%
    \pgfmathprintnumber[fixed,precision=2]{\pgfmathresult}\,\mbox{baselineskip}%
  \endgroup
}
\newcommand{\layoutGuideHorizontalBaselineDistance}[3][\campusheadertextgap]{%
  \begingroup
    \pgfextractx{\dimen0}{\pgfpointdiff{\pgfpointanchor{#2}{center}}{\pgfpointanchor{#3}{center}}}%
    \pgfmathparse{abs(\the\dimen0)/(1mm) / ((\the#1)/(1mm))}%
    \pgfmathprintnumber[fixed,precision=2]{\pgfmathresult}\,\mbox{baselineskip}%
  \endgroup
}
\begin{frame}{版式对齐检查}
  % Draw last, above Beamer's headline/footline, on this page only.
  \AddToHookNext{shipout/foreground}{%
    \begin{tikzpicture}[remember picture,overlay]
      \tikzset{layout guide/.style={red,line width=.4pt},
        layout label/.style={text=red,font=\tiny,fill=white,inner sep=.7pt},
        layout dimension/.style={layout guide,<->}}
      % Both transparent assets are square and cropped to the emblem bounds.
      \coordinate (emblem-left) at
        ([xshift=\dimexpr\campuslogoleftskip+\campuslogopadding\relax,
          yshift=\dimexpr-\campuslogoheight+\campuslogopadding+.5\campuslogoimagewidth\relax]
          current page.north west);
      \coordinate (emblem-right) at ([xshift=\campuslogoimagewidth]emblem-left);
      \coordinate (emblem-bottom) at
        ([xshift=\dimexpr\campuslogoleftskip+.5\campuslogowidth\relax,
          yshift=\dimexpr-\campuslogoheight+\campuslogopadding\relax]
          current page.north west);
      \draw[layout guide]
        ([yshift=.5\campuslogoimagewidth]emblem-left) rectangle
        ([yshift=-.5\campuslogoimagewidth]emblem-right);
      \draw[layout dimension] ([xshift=-\campuslogopadding]emblem-left) -- (emblem-left);
      \draw[layout dimension] (emblem-right) -- ([xshift=\campuslogopadding]emblem-right);
      \draw[layout dimension] (emblem-bottom) -- ([yshift=-\campuslogopadding]emblem-bottom);
      \node[layout label,anchor=north west] at
        ([xshift=\campuslogoleftskip,yshift=\dimexpr-\campuslogoheight-2mm\relax]
          current page.north west)
        {校徽左右/下方留白：\the\campuslogopadding};
      \foreach \offset/\tag in {
        \campuslogoleftskip/V1,
        {\dimexpr\campuslogoleftskip+\campuslogowidth\relax}/V2,
        {\dimexpr\campuslogoleftskip+\campusframetitleleftskip\relax}/V3,
        {\dimexpr\paperwidth-\campuslogoleftskip\relax}/V4}{%
        \coordinate (\tag) at ([xshift=\offset]current page.north west);
        \draw[layout guide] (\tag) -- (\tag |- current page.south);
        \node[layout label,anchor=north] at
          ([yshift=-16mm]\tag) {\tag};
      }
      \draw[layout guide,dashed] (current page.north) -- (current page.south);
      \node[layout label,anchor=north] at ([yshift=-16mm]current page.north) {中轴};

      % H1, H2 and H3 are the actual boundaries of the three header grid
      % rows: after the empty row, after the title row, and after the
      % subtitle row.
      \coordinate (H1) at ([yshift=-\campusheaderemptybaseline]current page.north west);
      \coordinate (H2) at ([yshift=-\dimexpr\campusheaderemptybaseline+\campusheaderbaseline\relax]current page.north west);
      \coordinate (H3) at ([yshift=-\dimexpr\campusheaderemptybaseline+\campusheaderbaseline+\campusheadersubtitlebaseline\relax]current page.north west);
      \coordinate (H5) at ([yshift=\dimexpr\campusnavigationfooterheight+\campusfootlineheight+\campusfootlinedepth\relax]current page.south west);
      \coordinate (H6) at ([yshift=\dimexpr\campussectionnavigationheight+\campusfootlineheight+2\campusfootlinedepth\relax]current page.south west);
      \coordinate (H7) at ([yshift=\dimexpr\campusfootlineheight+\campusfootlinedepth\relax]current page.south west);
      \coordinate (H8) at (current page.south west);
      % H4 is the body safety boundary: keep a baseline-grid breathing room
      % above the first footer band (H5).
      \coordinate (H4) at ([yshift=2.4\campusfooterbaseline]H5);
      \foreach \tag in {H1,H2,H3,H4,H5,H6,H7}{%
        \draw[layout guide] (\tag) -- (current page.east |- \tag);
        \node[layout label,anchor=east] at
          ([xshift=-2mm]current page.east |- \tag) {\tag};
      }
      \coordinate (safe-a) at ([xshift=15mm]current page.west |- H4);
      \coordinate (safe-b) at (safe-a |- H5);
      \draw[layout dimension] (safe-a) -- (safe-b)
        node[midway,right,layout label] {\layoutGuideBaselineDistance{safe-a}{safe-b}};
      % An inset half-stroke keeps the bottom boundary visible in the PDF.
      \draw[layout guide] ([yshift=.2pt]current page.south west)
        -- ([yshift=.2pt]current page.south east);

      % Show all three header grid rows separately.  These coordinates are
      % layout boundaries, so glyph ascenders and descenders do not change the
      % reported one-baseline lengths.
      \coordinate (header-top) at (current page.north west);
      \foreach \upper/\lower/\pos/\label/\unit in {
        header-top/H1/.70/上方空行 (\string\tiny)/\campusheaderemptybaseline,
        H1/H2/.78/标题 (\string\Large)/\campusheaderbaseline,
        H2/H3/.86/副标题 (\string\normalsize)/\campusheadersubtitlebaseline}{%
        \coordinate (header-a) at ([xshift=\pos\paperwidth]current page.west |- \upper);
        \coordinate (header-b) at (header-a |- \lower);
        \draw[layout dimension] (header-a) -- (header-b)
          node[midway,left,layout label] {\label：\layoutGuideBaselineDistance[\unit]{header-a}{header-b}};
      }

      % Horizontal dimensions occupy the open area below the diagnostic text.
      \coordinate (logo-a) at ([yshift=-63mm]V1);
      \coordinate (logo-b) at (V2 |- logo-a);
      \draw[layout dimension] (logo-a) -- (logo-b)
        node[midway,above,anchor=south west,layout label] {校徽色块：%
          \layoutGuideHorizontalBaselineDistance[\campuslogowidth]{logo-a}{logo-b}
          (\string\Huge)};
      \coordinate (gap-a) at ([yshift=-67mm]V2);
      \coordinate (gap-b) at (V3 |- gap-a);
      \draw[layout dimension] (gap-a) -- (gap-b)
        node[midway,above,anchor=south west,layout label] {%
          \layoutGuideHorizontalBaselineDistance{gap-a}{gap-b} (\string\normalsize)};
      \coordinate (body-a) at ([yshift=-70mm]V1);
      \coordinate (body-b) at (V4 |- body-a);
      \draw[layout dimension] (body-a) -- (body-b)
        node[midway,above,layout label] {正文宽度：\layoutGuideDistance{body-a}{body-b}};
      \coordinate (margin-a) at ([yshift=-72mm]current page.north west);
      \coordinate (margin-b) at (V1 |- margin-a);
      \draw[layout dimension] (margin-a) -- (margin-b)
        node[midway,above,anchor=south west,layout label] {%
          \layoutGuideHorizontalBaselineDistance[\campusheaderbaseline]{margin-a}{margin-b} (\string\Large)};
      \coordinate (margin-c) at (V4 |- margin-a);
      \coordinate (margin-d) at (current page.east |- margin-a);
      \draw[layout dimension] (margin-c) -- (margin-d)
        node[midway,above,anchor=south east,layout label] {%
          \layoutGuideHorizontalBaselineDistance[\campusheaderbaseline]{margin-c}{margin-d} (\string\Large)};

      % Footer heights are read from the theme, so changes remain in sync.
      % Callouts above the bands keep their text clear of the navigation.
      \foreach \upper/\lower/\pos/\label in {
        H5/H6/.55/subsection (\string\scriptsize),
        H6/H7/.72/section (\string\scriptsize),
        H7/H8/.76/信息行 (\string\scriptsize)}{%
        \coordinate (band-a) at ([xshift=\pos\paperwidth]current page.west |- \upper);
        \coordinate (band-b) at (band-a |- \lower);
        \draw[layout dimension] (band-a) -- (band-b)
          node[midway,right,layout label] {\label：\layoutGuideBaselineDistance{band-a}{band-b}};
      }
    \end{tikzpicture}%
  }%
  \vspace*{11mm}
  \hspace*{-2pt}%
  \begin{minipage}{\dimexpr\textwidth-6pt\relax}
    \begin{description}[竖线]
      \item[竖线] V1 正文左界；V2 校徽色块右界；V3 页眉文字起点；V4 正文右界。
      \item[页眉] H1/H2/H3 为三行网格分界线；箭头以对应字号的 baselineskip 为单位。
      \item[页脚] H5--H7 为三行色带的上边界；间距以 \texttt{\string\scriptsize} 字号的 baselineskip 为单位。
      \item[参考] H4 为正文安全线，安全余量按页脚基线网格设置；虚线为页面中轴。
    \end{description}
  \end{minipage}
\end{frame}
\note{红线用于核对校徽、标题、正文和页脚的对齐。正式汇报通常不保留这张诊断页。}
\endgroup

% Read sizes from the active LaTeX commands rather than duplicating pt values.
\begingroup
\makeatletter
\newcommand{\fontGuideSize}[1]{%
  \begingroup #1\edef\fontGuideValue{\f@size}%
    \footnotesize\fontGuideValue\,pt\endgroup}
\makeatother
\newcommand{\fontGuideSample}[1]{%
  \makebox[24mm][l]{{\color{maincolor}#1 字 Aa}}}
\begin{frame}{模板字号与用途}
  \footnotesize
  \setlength{\tabcolsep}{0pt}
  \renewcommand{\arraystretch}{1.05}
  \begin{tabular}{@{}l@{\hspace{2mm}}p{28mm}@{\hspace{2mm}}p{15mm}@{\hspace{2mm}}p{72mm}@{}}
    \textbf{实际大小} & \textbf{字号命令} & \textbf{字号} & \textbf{主要用途} \\[.15mm]
    \fontGuideSample{\Huge\bfseries} & \texttt{\string\Huge} & \fontGuideSize{\Huge} &
      封面/章节/致谢页主标题 \\[.15mm]
    \fontGuideSample{\LARGE} & \texttt{\string\LARGE} & \fontGuideSize{\LARGE} &
      单行页眉（目录/参考文献等）；致谢页两行感谢语 \\[.15mm]
    \fontGuideSample{\Large\bfseries} & \texttt{\string\Large} & \fontGuideSize{\Large} &
      双行页眉上行、section 编号；封面／致谢页副标题 \\[.15mm]
    \fontGuideSample{\large} & \texttt{\string\large} & \fontGuideSize{\large} &
      可选的正文强调字号 \\[.15mm]
    \fontGuideSample{\normalsize} & \texttt{\string\normalsize} & \fontGuideSize{\normalsize} &
      正文/页眉下行；课题组/个人信息；小节导读 \\[.15mm]
    \fontGuideSample{\small} & \texttt{\string\small} & \fontGuideSize{\small} &
      提示框正文、二级列表；本例参考文献 \\[.15mm]
    \fontGuideSample{\footnotesize} & \texttt{\string\footnotesize} & \fontGuideSize{\footnotesize} &
      提示框标题、三级列表、图表题注 \\[.15mm]
    \fontGuideSample{\scriptsize} & \texttt{\string\scriptsize} & \fontGuideSize{\scriptsize} &
      三行页脚（两行导航与信息行） \\[.15mm]
    \fontGuideSample{\tiny} & \texttt{\string\tiny} & \fontGuideSize{\tiny} &
      红色标注、角落引用、论文元信息
  \end{tabular}
  \par\smallskip
  {\scriptsize 字号由当前文档实际读取，单位为 TeX pt。粗细由元素样式决定；数学上下标自动缩小。\par
  参考文献字号可通过命令参数调整；使用 shrink 时，整页缩放可能使最终显示更小。\par}
\end{frame}
\note{比较不同字号的实际显示。正文以正常字号为主，代码使用 small；内容过多时优先拆页。}
\endgroup


% section 组织整节内容，subsection 自动显示在页眉上方。
% frame 的标题参数即 frametitle，只填写本页标题，无需额外设置副标题。
\begin{frame}{三级内容结构}
      \begin{description}[Subsection]
            \item[Section] 组织整节内容，用于起始页和页脚导航。
            \item[Subsection] 组织节内主题，自动显示在页眉上方。
            \item[Frametitle] 概括本页内容，在页眉下方以小字号显示。
            \item[连续页面] 同一 subsection 下可放多页，每页只需填写自己的标题。
      \end{description}
\end{frame}
\note{说明 section 管整节，subsection 显示在页眉上行，frametitle 概括当前页。转入具体内容示例。}

% 公式与图像沿用本 subsection，保留单小节的导航示例。
% Formula examples use the amsmath environments already loaded by Beamer.
\begin{frame}{行内公式与多行推导}
  对样本 $\mathbf{x}_i\in\mathbb{R}^{d}$，用线性模型预测标量目标 $y_i$：
  \begin{align*}
    \widehat{y}_i &= \mathbf{w}^{\mathsf T}\mathbf{x}_i+b, \\
    \mathcal{L}(\mathbf{w},b)
      &= \frac{1}{N}\sum_{i=1}^{N}(\widehat{y}_i-y_i)^2
         +\lambda\lVert\mathbf{w}\rVert_2^2, \\
    \nabla_{\mathbf{w}}\mathcal{L}
      &= \frac{2}{N}\sum_{i=1}^{N}(\widehat{y}_i-y_i)\mathbf{x}_i
         +2\lambda\mathbf{w}.
  \end{align*}
  \begin{notebox}[排版要点]
    行内公式用 \texttt{\$...\$}；多行公式用 \texttt{align*}，
    在等号前放置 \texttt{\&} 对齐，星号表示不显示编号。
  \end{notebox}
\end{frame}
\note{沿等号解释预测、损失与梯度，强调多行推导应对齐关键运算符。这是数学排版示例。}

\begin{frame}{矩阵、分段函数与公式说明}
  \begin{columns}[T,onlytextwidth]
    \begin{column}{.47\textwidth}
      \textbf{矩阵：批量预测}
      \[
        \mathbf{X}=\begin{bmatrix}
          \mathbf{x}_1^{\mathsf T}\\
          \vdots\\
          \mathbf{x}_N^{\mathsf T}
        \end{bmatrix},\qquad
        \widehat{\mathbf{y}}=\mathbf{X}\mathbf{w}+b\mathbf{1}.
      \]
      \texttt{bmatrix} 自动生成方括号；矩阵中的列用 \texttt{\&} 分隔。
    \end{column}
    \begin{column}{.47\textwidth}
      \textbf{分段函数：ReLU}
      \[
        f(z)=\begin{cases}
          z, & z\geq 0,\\
          0, & z<0.
        \end{cases}
      \]
      \texttt{cases} 排列函数值与条件；公式中的文字可用 \texttt{\string\text\{...\}}。
    \end{column}
  \end{columns}
  \medskip
  \begin{tipbox}[符号与说明]
    $N$ 是样本数，$d$ 是特征维数；$\mathbf{1}\in\mathbb{R}^{N}$ 为全一向量。
    公式与符号解释放在同一页，便于阅读。
  \end{tipbox}
\end{frame}
\note{左侧检查矩阵间距，右侧检查分段条件。下方强调符号定义应在使用前说明。}


% !TeX root = ../example.tex
% Keep these examples within the existing subsection/navigation structure.
\begin{frame}[fragile]{代码展示：Python}
\begin{campuscode}[language=Python,numbers,caption={保留缩进，便于讲解函数结构。}]{normalize.py}
def normalize(values):
    # 将数值归一化，保留清晰的缩进
    total = sum(values)
    if total == 0:
        return [0.0 for value in values]
    return [value / total for value in values]

print(normalize([2, 3, 5]))
\end{campuscode}
\small 标题栏标识文件与语言；下方说明补充阅读提示。
\end{frame}

\begin{frame}[fragile]{代码展示：终端命令与输出}
\begin{campuscode}[language=bash,caption={静态展示命令与输出，不在编译时执行。}]{Terminal}
$ python -m venv .venv
$ source .venv/bin/activate
$ python normalize.py
[0.2, 0.3, 0.5]
\end{campuscode}
\small 用命令提示符区分输入与输出。
\end{frame}

\begin{frame}[fragile]{代码展示：从文件读取}
\campusinputcode[language=Python,description={从文件读取并高亮选定源码。}]
  {normalize.py}{chapters/code/normalize.py}
\small 用 \verb|\campusinputcode[language=Python]{文件名}{文件路径}| 复用源文件。
\end{frame}


\begin{frame}[fragile]{自动适配素材}
      \begin{columns}[T,onlytextwidth]
            \begin{column}{.48\textwidth}
                  \textbf{适配素材并添加图注}
                  \begin{verbatim}
\begin{figure}
  \fitgraphic[
    width=.42,
    height=.40
  ]{assets/sigs_wordmark.png}
  \caption{标志素材的适配效果}
\end{figure}
                  \end{verbatim}
                  宽、高分别是整页宽、高的比例，图片在范围内保持纵横比。
            \end{column}
            \begin{column}{.48\textwidth}
                  \begin{figure}
                        \fitgraphic[width=.42,height=.40]{assets/sigs_wordmark.png}
                        \caption{标志素材的适配效果}
                  \end{figure}
            \end{column}
      \end{columns}
\end{frame}
\note{对照左右两栏：左侧是插图源码，右侧是实际输出。宽高参数表示整页尺寸的比例，图片保持纵横比。}

\begin{frame}[fragile]{自动适配图片与图注}
      \begin{columns}[T,onlytextwidth]
            \begin{column}{.48\textwidth}
                  \textbf{同时设置图像与图注}
                  \begin{verbatim}
\fitfigure[
  width=.42,
  height=.28
]{\schoolwordmarkonlight}
 {横向标志素材的适配效果}
                  \end{verbatim}
                  最后一项填写图注，留空则不显示。
            \end{column}
            \begin{column}{.48\textwidth}
                  \fitfigure[width=.42,height=.28]{\schoolwordmarkonlight}%
                        {横向标志素材的适配效果}
            \end{column}
      \end{columns}
\end{frame}
\note{fitfigure 同时提供适配与图注。最后一个参数留空时不显示图注，便于简洁展示。}

\begin{rightimageframe}[image width=.40]{右侧满高图片}{assets/tsinghua_temple.png}
      \begin{description}[高度]
            \item[高度] 图片从页面顶端延伸到导航上边缘。
            \item[宽度] \texttt{image width} 控制右侧图片区域占页面宽度的比例。
            \item[正文] 左侧仍可放置列表、公式和提示框。
      \end{description}
\end{rightimageframe}
\note{检查右侧图片从页顶延伸至导航上沿，正文留在左侧。图片路径由当前汇报指定。}

\begin{frame}[fragile]{图像命令速查}
      \begin{description}[适配素材]
            \item[适配素材] \verb|\fitgraphic[width=.8,height=.6]{file}|
            \item[添加图注] \verb|\fitfigure[width=.8,height=.6]{file}{caption}|
            \item[右侧图片] \verb|\begin{rightimageframe}{标题}{file}|
            \item[PDF 页面] 可选参数 \texttt{page=2} 指定 PDF 的第 2 页。
      \end{description}
\end{frame}
\note{回顾三个常用图像入口；多页 PDF 可通过 page 参数选页。之后转入论文引用示例。}
